\newpage
\section{Cel ćwiczenia}

Celem zajęć było złożenie układu kombinacyjnego od opisu słownego aż do programu w sterowniku. Dla trzech zadań przeszliśmy tę samą drogę: opis, wejścia i wyjścia, tablica prawdy, funkcja w postaci kanonicznej, minimalizacja, schemat z bramek i drabinka w LOGO! Soft Comfort. Testu na zaciskach w tym sprawozdaniu nie ma. Poprawność sprawdzaliśmy w symulatorze, na stanach, które rozdzielają gałęzie funkcji.

Układ kombinacyjny liczy wyjście wyłącznie z tego, co w danym obiegu jest na wejściach. Nie wstawialiśmy markerów ani timerów, bo pamięć poprzedniego stanu zmieniłaby zadanie w układ sekwencyjny.

Po ćwiczeniu potrafimy:

\begin{itemize}
    \item rozpisać wszystkie kombinacje wejść i odczytać z nich sumę iloczynów,
    \item zwinąć tę sumę do postaci, która mieści się na kilku stykach,
    \item narysować to samo bramkami AND, OR i NOT,
    \item przenieść bramki na drabinkę: szereg jako iloczyn, gałąź równoległą jako sumę, styk zanegowany jako negację,
    \item sprawdzić w symulacji, która gałąź faktycznie trzyma cewkę.
\end{itemize}

\section{Stanowisko}

Program rysowaliśmy w LOGO! Soft Comfort jako schemat drabinkowy i odpalaliśmy go w symulacji wbudowanej w to samo środowisko. Sterownik LOGO! czyta wejścia, wykonuje program i wystawia wyjścia, a potem powtarza obieg. Przy funkcji kombinacyjnej jeden taki obieg wystarcza: cewka nie zależy od tego, co było na niej cykl wcześniej.

Wejścia w LOGO! numerują się od I1, wyjścia od Q1. W instrukcji ćwiczenia te same sygnały są zapisane adresami typu I0.0, tak jak w sterownikach S7. Zostawiliśmy numerację z programu, bo tak podpisane są bloki na zrzutach.

W symulacji przełącznik wejścia odsunięty w stronę zera jest niebieski, odsunięty w stronę jedynki jest czerwony. Żarówka przy Q1 świeci na żółto, gdy cewka jest załączona, i gaśnie, gdy cewka odpadła.
